\label{chap:fm3-planted-forest-synthesis}

\providecommand{\MC}{\mathrm{MC}}
\providecommand{\Defcyc}{\operatorname{Def}_{\mathrm{cyc}}}
\providecommand{\Cone}{\operatorname{Cone}}
\providecommand{\Tot}{\operatorname{Tot}}
\providecommand{\adop}{\operatorname{ad}}
\providecommand{\cA}{\mathcal A}
\providecommand{\cB}{\mathcal B}
\providecommand{\cC}{\mathcal C}
\providecommand{\cD}{\mathcal D}
\providecommand{\cF}{\mathcal F}
\providecommand{\cL}{\mathcal L}
\providecommand{\cR}{\mathcal R}
\providecommand{\Q}{\mathbf Q}
\providecommand{\C}{\mathbf C}
\providecommand{\Pone}{\mathbf P^1}
\providecommand{\Aone}{\mathbf A^1}

\title{Modular homotopy theory from first principles:\\
residue splitting, planted forests, and the explicit \texorpdfstring{$\FM_3(X)$}{FM3(X)} model}
\author{}
\date{}

\begin{abstract}
This note synthesizes a complete thread of ideas into a single mathematical arc.
One begins with the modular Maurer--Cartan package for chiral Koszul duality,
passes to the renormalization problem on Fulton--MacPherson compactifications,
splits the logarithmic deformation algebra into residue and counterterm pieces,
and then compresses the first nontrivial case $\FM_3(X)$ to an explicit small
$L_\infty$ model.  The local counterterm sector is contractible; the surviving
local theory is residue-theoretic; and the first genuinely higher bracket is a
nested-collision operation supported on the codimension-two planted forests
$ij\prec 123$.  The note also records the string-theoretic blueprint that this
geometry suggests: worldsheet sewing becomes planted-forest gluing, the open
sector lives on the bar side, the closed sector on the cobar or defect-cyclic
side, branes force ordered $E_1$-behaviour, and the first planted-forest cubic
obstruction is the smallest local shadow of higher-genus holographic transport.
\end{abstract}

\tableofcontents

\section{The governing mechanism}

The heart of the subject is a single piece of structure: a chiral algebra
$\cA$ on a curve carries a bar object whose total genus-corrected differential
packages boundary collisions, modular sewing, and deformation theory at once.
At the model level one obtains a complete filtered dg Lie algebra of positive-genus
coderivations, and the total differential determines a canonical Maurer--Cartan
element there.  The global cyclic deformation package is the abstract target;
Fulton--MacPherson geometry is the mechanism that makes the package computable.

The conceptual chain is:
\begin{equation}
\label{eq:grand-chain}
\text{bar/coderivation geometry}
\Longrightarrow
\text{cyclic deformation }L_\infty\text{ algebra}
\Longrightarrow
\text{local logarithmic renormalization}
\Longrightarrow
\text{residue small models on planted forests}.
\end{equation}
The point of the present note is to make the last implication explicit.

\section{The coderivation Maurer--Cartan envelope}

Let $\bar B_X^{\mathrm{full}}(\cA)=\prod_{g\ge 0}\hbar^g\bar B_X^{(g)}(\cA)$ be the
completed genus-graded bar object, equipped with a total differential
\[
D = D^{(0)} + \delta_{\cA},
\qquad
D^2=0,
\qquad
(D^{(0)})^2=0.
\]
Let $F^p$ denote the positive-genus filtration.

\begin{definition}[Modular coderivation envelope]
The positive-genus coderivation dg Lie algebra is
\[
\mathfrak g^{\mathrm{mod}}_{\mathrm{cod}}(\cA)
:=
\bigl(F^1\operatorname{Coder}(\bar B_X^{\mathrm{full}}(\cA)),\; \partial=[D^{(0)},-],\;[-,-]\bigr).
\]
\end{definition}

\begin{theorem}[Canonical coderivation Maurer--Cartan class; \ClaimStatusProvedHere]
\label{thm:coderivation-MC-synthesis}
The element
\[
\delta_{\cA}:=D-D^{(0)}
\in
F^1\operatorname{Coder}(\bar B_X^{\mathrm{full}}(\cA))
\]
is a Maurer--Cartan element of $\mathfrak g^{\mathrm{mod}}_{\mathrm{cod}}(\cA)$:
\[
\partial\delta_{\cA}+\frac12[\delta_{\cA},\delta_{\cA}] = 0.
\]
\end{theorem}

\begin{proof}
Expand $D^2=0$ using $D=D^{(0)}+\delta_{\cA}$ and $(D^{(0)})^2=0$:
\[
0 = [D^{(0)},\delta_{\cA}] + \delta_{\cA}^2
= \partial\delta_{\cA}+\frac12[\delta_{\cA},\delta_{\cA}].
\]
The filtration is pronilpotent because the bracket of a genus-$p$ coderivation
with a genus-$q$ coderivation has genus at least $p+q$.
\end{proof}

This shifts the nonlinear problem from ``construct a modular Maurer--Cartan
class from scratch'' to ``compare the coderivation class to the desired cyclic
model''.  The local geometry below should be read as the smallest explicit model
for that comparison problem.

\section{Filtered transfer and genus-by-genus obstruction theory}

Let $L$ be a complete filtered cyclic $L_\infty$ algebra, and suppose there is a
filtered $L_\infty$ quasi-isomorphism
\[
\Phi:\mathfrak g^{\mathrm{mod}}_{\mathrm{cod}}(\cA)
\rightsquigarrow
L.
\]
Then the image of $\delta_{\cA}$ is a Maurer--Cartan element of $L$.
This is formal $L_\infty$ functoriality.  The non-formal question is how one
actually computes the image.  Theorem~\ref{thm:fm3-higher-obstruction-synthesis}
below answers that in the smallest boundary case.

Independently of geometry, if one seeks
\[
\Theta = \sum_{g\ge 1}\theta_g
\]
with $\theta_g$ in filtration $g$, the Maurer--Cartan equation expands
recursively.  The genus-$g$ obstruction is always a degree-$3$ cohomology class,
and rigidity is controlled by degree-$1$ and degree-$2$ cohomology.  The local
Fulton--MacPherson theory constructed below makes these classes explicit in terms
of residues on boundary strata.

\section{Local logarithmic renormalization from first principles}

Let $(Y,\partial Y)$ be a smooth simple normal crossings pair.  Near a stratum
\(S=\{x_1=\cdots=x_r=0\}\) one has logarithmic forms generated by
\[
\eta_i = d\log x_i.
\]
Let $\mathfrak g$ be a complete filtered dg Lie algebra.  The completed
logarithmic deformation algebra is
\[
\widehat\Omega^\bullet_{\widehat Y_S}(\log \partial Y)\widehat\otimes \mathfrak g.
\]
The ideal of the stratum is $J_S=(x_1,\dots,x_r)$.

\begin{definition}[Local counterterm sector and residue quotient]
Define
\[
\widehat{\mathcal I}_S(\mathfrak g)
:=
J_S\cdot \widehat\Omega^\bullet_{\widehat Y_S}(\log \partial Y)\widehat\otimes \mathfrak g,
\qquad
\mathcal R_S(\mathfrak g)
:=
\widehat\Omega^\bullet_{\widehat Y_S}(\log \partial Y)\widehat\otimes \mathfrak g\,/\,\widehat{\mathcal I}_S(\mathfrak g).
\]
\end{definition}

\begin{theorem}[Residue splitting theorem; \ClaimStatusProvedHere]
\label{thm:local-residue-splitting-synthesis}
The ideal $\widehat{\mathcal I}_S(\mathfrak g)$ is a contractible dg Lie ideal, and
\[
\mathcal R_S(\mathfrak g)
\cong
\Omega^\bullet_S(\log \partial S)\widehat\otimes
\Lambda^\bullet N^{\log,\vee}_{S/Y}\widehat\otimes \mathfrak g.
\]
Consequently the projection to the residue quotient is a filtered dg Lie
quasi-isomorphism.
\end{theorem}

\begin{proof}
Let
\[
E=\sum_{i=1}^r x_i\frac{\partial}{\partial x_i}
\]
be the normal Euler field.  On a monomial $x^\alpha\eta_I\wedge dy_J$ of positive
normal degree $|\alpha|$ one has
\[
\mathcal L_E(x^\alpha\eta_I\wedge dy_J)=|\alpha|\,x^\alpha\eta_I\wedge dy_J.
\]
Define
\[
h = \frac1{|\alpha|}\iota_E
\]
on the positive-degree part.  Cartan's formula gives
\[
(dh+hd)\omega = \omega
\]
for every element of positive normal degree.  Thus the counterterm sector is
contractible.  Modding out by positive normal degree leaves precisely the residue
part generated by forms on $S$ and the logarithmic normal classes $\eta_i$.
\end{proof}

This theorem says: the local singular part is homotopically trivial; the only
nontrivial local data is residue data.  The smallest case in which that residue
data already produces a genuinely higher bracket is $\FM_3(X)$.

\section{The global planted-forest principle}

On a Fulton--MacPherson compactification $\FM_n(X)$ or its higher-genus analogues,
the boundary is indexed by collision forests.  After local residue reduction, the
global logarithmic deformation algebra becomes a semicosimplicial residue dg Lie
algebra indexed by those forests.  Compressing this semicosimplicial object to a
small model transfers the strict dg Lie structure to an $L_\infty$ structure
whose higher operations record precisely the failure of local residue data to
match globally.

The smallest nontrivial instance is $\FM_3(X)$, because its planted forests are
\[
12,\ 23,\ 13,\ 123,\ 12\prec 123,\ 23\prec 123,\ 13\prec 123.
\]
There is no smaller space on which a nested collision exists.

\section{The boundary geometry of \texorpdfstring{$\FM_3(X)$}{FM3(X)}}

Let $X$ be a smooth complex curve.
The compactification $\FM_3(X)$ is obtained by blowing up the small diagonal and
then the proper transforms of the three pair diagonals.  The boundary divisors
are
\[
D_{12},\ D_{23},\ D_{13},\ D_{123},
\]
where $D_{123}$ is the exceptional divisor over the small diagonal.  The three
pair divisors are disjoint, each meets $D_{123}$ transversely, and the only
codimension-two strata are
\[
S_{12}=D_{12}\cap D_{123},\qquad
S_{23}=D_{23}\cap D_{123},\qquad
S_{13}=D_{13}\cap D_{123}.
\]
Moreover,
\[
D_{ij}\cong X^2,
\qquad
D_{123}\to X \text{ is a }\Pone\text{-bundle},
\qquad
S_{ij}\cong X.
\]

A local coordinate proof is elementary.  Near the small diagonal in $X^3$, use
\[
(x,u,v)=(z_1,z_2-z_1,z_3-z_1).
\]
Then the pair diagonals are the three lines $u=0$, $v=0$, $u-v=0$ in the normal
plane.  After blowing up the origin, their strict transforms are disjoint and
meet the exceptional $\Pone$ at the three points $[0:1],[1:0],[1:1]$.

\section{The residue dg Lie algebras of the boundary strata}

Fix a complete filtered dg Lie algebra $\mathfrak g$.  Introduce logarithmic
normal classes
\[
\eta_{12},\eta_{23},\eta_{13},\eta_{123}
\]
along the four boundary divisors.  Define the divisor residue dg Lie algebras
\begin{align*}
\mathfrak r_{12} &:= \Omega^\bullet(D_{12})\widehat\otimes \Lambda\langle\eta_{12}\rangle\widehat\otimes \mathfrak g,\\
\mathfrak r_{23} &:= \Omega^\bullet(D_{23})\widehat\otimes \Lambda\langle\eta_{23}\rangle\widehat\otimes \mathfrak g,\\
\mathfrak r_{13} &:= \Omega^\bullet(D_{13})\widehat\otimes \Lambda\langle\eta_{13}\rangle\widehat\otimes \mathfrak g,\\
\mathfrak r_{123} &:= \Omega^\bullet(D_{123})\widehat\otimes \Lambda\langle\eta_{123}\rangle\widehat\otimes \mathfrak g,
\end{align*}
and the codimension-two residue dg Lie algebras
\begin{align*}
\mathfrak s_{12} &:= \Omega^\bullet(S_{12})\widehat\otimes \Lambda\langle\eta_{12},\eta_{123}\rangle\widehat\otimes \mathfrak g,\\
\mathfrak s_{23} &:= \Omega^\bullet(S_{23})\widehat\otimes \Lambda\langle\eta_{23},\eta_{123}\rangle\widehat\otimes \mathfrak g,\\
\mathfrak s_{13} &:= \Omega^\bullet(S_{13})\widehat\otimes \Lambda\langle\eta_{13},\eta_{123}\rangle\widehat\otimes \mathfrak g.
\end{align*}
Set
\[
\mathfrak r_0:=\mathfrak r_{12}\oplus\mathfrak r_{23}\oplus\mathfrak r_{13}\oplus\mathfrak r_{123},
\qquad
\mathfrak r_1:=\mathfrak s_{12}\oplus\mathfrak s_{23}\oplus\mathfrak s_{13}.
\]
The planted-forest incidence morphism is
\[
\chi(\alpha_{12},\alpha_{23},\alpha_{13},\alpha_{123})
=
(\alpha_{12}|_{S_{12}}-\alpha_{123}|_{S_{12}},\
 \alpha_{23}|_{S_{23}}-\alpha_{123}|_{S_{23}},\
 \alpha_{13}|_{S_{13}}-\alpha_{123}|_{S_{13}}).
\]
It is a filtered dg Lie morphism.

\section{The explicit planted-forest \texorpdfstring{$L_\infty$}{L-infinity} model}

Define the shifted mapping cone
\[
\mathfrak P_{\FM_3}(X,\mathfrak g):=\Cone(\chi)[-1].
\]
An element of degree $k$ is a pair
\[
(\alpha,\beta),\qquad \alpha\in \mathfrak r_0^k,\quad \beta\in \mathfrak r_1^{k-1}.
\]

\begin{theorem}[The explicit small model; \ClaimStatusProvedHere]
\label{thm:fm3-small-model-synthesis}
The boundary logarithmic deformation algebra of $\FM_3(X)$ is filtered
$L_\infty$ quasi-isomorphic to $\mathfrak P_{\FM_3}(X,\mathfrak g)$.  The first
three brackets are:
\begin{align}
\label{eq:syn-l1}
l_1(\alpha,\beta) &= (d\alpha,\chi(\alpha)-d\beta),\\
\label{eq:syn-l2}
l_2((\alpha,\beta),(\alpha',\beta'))
&=
\left([
\alpha,\alpha'],\ \frac12\bigl([\beta,\chi(\alpha')] - (-1)^{|x||x'|}[\beta',\chi(\alpha)]\bigr)\right),\\
\label{eq:syn-l3}
l_3((0,\beta_1),(0,\beta_2),(\alpha,0))
&=
\frac1{12}\Bigl(0,
[\beta_1,[\beta_2,\chi(\alpha)]]+(-1)^{|m_1||m_2|}[\beta_2,[\beta_1,\chi(\alpha)]]\Bigr),
\end{align}
with all other values obtained by graded skew-symmetry or vanishing.  In
particular, the first genuinely higher operation is supported exactly on the
nested forests $ij\prec 123$.
\end{theorem}

\begin{proof}
The boundary cover of $\FM_3(X)$ has four patches and only three nonempty
pairwise overlaps.  After residue reduction, its semicosimplicial dg Lie algebra
is therefore a two-term diagram $\mathfrak r_0\to\mathfrak r_1$.
The Thom--Whitney totalization of this diagram is a strict dg Lie algebra, and
its normalized cone inherits the displayed $L_\infty$ operations by homotopy
transfer.  The coefficients are read off from the Maurer--Cartan equation of the
homotopy fiber, expanded to cubic order.  The only possible cubic terms come
from one divisor input and two overlap inputs, because there are no triple
overlaps in the cover.
\end{proof}

\section{Closed-form Maurer--Cartan theory}

\begin{theorem}[Maurer--Cartan dictionary; \ClaimStatusProvedHere]
\label{thm:fm3-mc-synthesis}
A degree-one element $(\alpha,\beta)\in \mathfrak r_0^1\oplus\mathfrak r_1^0$ is
Maurer--Cartan in $\mathfrak P_{\FM_3}(X,\mathfrak g)$ if and only if
\begin{align}
\label{eq:syn-alpha-mc}
d\alpha + \frac12[\alpha,\alpha] &= 0,\\
\label{eq:syn-beta-mc}
d\beta &= \frac{\adop_\beta}{1-e^{-\adop_\beta}}\chi(\alpha)
= \chi(\alpha)+\frac12[\beta,\chi(\alpha)] + \frac1{12}[\beta,[\beta,\chi(\alpha)]] - \cdots.
\end{align}
Equivalently,
\[
e^\beta * \chi(\alpha)=0.
\]
\end{theorem}

\begin{proof}
The homotopy fiber of a dg Lie morphism classifies Maurer--Cartan elements on
the source together with a gauge trivialization of their image in the target.
Applying this to $\chi$ gives the claim.  The Bernoulli expansion is the power
series identity
\[
\frac{z}{1-e^{-z}}=1+\frac z2+\frac{z^2}{12}-\frac{z^4}{720}+\cdots.
\]
\end{proof}

This formula is the most efficient way to understand the geometry:
$\alpha$ is divisor residue data, $\chi(\alpha)$ is the planted-forest mismatch
along the codimension-two strata, and $\beta$ is the overlap parameter that
tries to gauge that mismatch away.

\section{The first genuinely planted-forest obstruction}

Let
\[
\Theta(t)=t\Theta_1+t^2\Theta_2+t^3\Theta_3+\cdots,
\qquad
\Theta_k=(\alpha_k,\beta_k).
\]

\begin{proposition}[First and second obstruction classes; \ClaimStatusProvedHere]
\label{prop:fm3-first-second-synthesis}
The order-$t$ equations are
\[
d\alpha_1=0,
\qquad
 d\beta_1=\chi(\alpha_1).
\]
Hence the first gluing obstruction is
\[
o_1(\alpha_1)=[\chi(\alpha_1)]\in H^2(\mathfrak P_{\FM_3},l_1)\cong H^1(\mathfrak r_1,d).
\]
If this vanishes and $\Theta_1$ is chosen, the quadratic obstruction is
\[
o_2(\Theta_1)=\frac12l_2(\Theta_1,\Theta_1)
=\left(\frac12[\alpha_1,\alpha_1],\frac12[\beta_1,\chi(\alpha_1)]\right)
\in Z^2(\mathfrak P_{\FM_3},l_1).
\]
\end{proposition}

\begin{theorem}[The first genuinely planted-forest cubic obstruction; \ClaimStatusProvedHere]
\label{thm:fm3-higher-obstruction-synthesis}
Assume $\Theta_1$ and $\Theta_2$ solve the Maurer--Cartan equation modulo $t^3$.
Then the obstruction to extending to order $t^3$ is
\[
o^{\mathrm{pf}}_3(\Theta_1,\Theta_2)
=
\left[l_2(\Theta_1,\Theta_2)+\frac16l_3(\Theta_1,\Theta_1,\Theta_1)\right]
\in H^2(\mathfrak P_{\FM_3},l_1).
\]
Its overlap component is
\begin{equation}
\label{eq:pf-cubic-overlap-synthesis}
\frac12\bigl([\beta_1,\chi(\alpha_2)] + [\beta_2,\chi(\alpha_1)]\bigr)
+ \frac1{12}[\beta_1,[\beta_1,\chi(\alpha_1)]].
\end{equation}
Componentwise, on the nested forest $ij\prec 123$ this is
\[
\frac12\bigl([\beta_{1,ij},\chi_{ij}(\alpha_2)] + [\beta_{2,ij},\chi_{ij}(\alpha_1)]\bigr)
+ \frac1{12}[\beta_{1,ij},[\beta_{1,ij},\chi_{ij}(\alpha_1)]].
\]
\end{theorem}

\begin{proof}
This is the $t^3$ coefficient of the Maurer--Cartan equation using
\eqref{eq:syn-l2} and \eqref{eq:syn-l3}.  The first two terms come from the
binary bracket; the last term is the cubic planted-forest correction coming from
nested collisions.
\end{proof}

The crucial feature of \eqref{eq:pf-cubic-overlap-synthesis} is that it has no
analogue on $\FM_2(X)$.  It is the first place where the codimension-two
stratification contributes a genuinely higher local correction.

\section{Specialization to the cyclic deformation algebra}

Taking $\mathfrak g=\Defcyc(\cA)$ gives a local small model for the boundary part
of the cyclic deformation algebra of a chiral algebra $\cA$.
The divisor variables $\alpha_{ij}$ and $\alpha_{123}$ are residue incarnations
of the modular deformation data; the overlap variables $\beta_{ij}$ measure the
planted-forest gauge glue between the pair-collision residues and the triple
residue.  The first nontrivial local higher operation is the cubic defect
\eqref{eq:pf-cubic-overlap-synthesis}.

\section{String-theoretic blueprint}

The entire picture has a direct string-theoretic reading.

\subsection*{Worldsheet}
Fulton--MacPherson blowups are the local algebraic avatars of worldsheet sewing.
The planted forests of $\FM_3(X)$ are the first nontrivial sewing types.  The
incidence map $\chi$ is the algebraic sewing defect, and the cubic obstruction
of Theorem~\ref{thm:fm3-higher-obstruction-synthesis} is the first local
higher-collision correction.

\subsection*{Open strings}
On the boundary side, the bar construction is the natural home of the open
BRST complex.  The overlap parameter $\beta$ in the $\FM_3$ model is the first
local gluing field needed to make boundary residues compatible across nested
collisions.

\subsection*{Closed strings}
The closed sector should be recovered from a cobar or defect-cyclic bulk object.
The equation
\[
d\beta = \frac{\adop_\beta}{1-e^{-\adop_\beta}}\chi(\alpha)
\]
identifies the precise local condition that any such bulk object must satisfy in
order to absorb nested boundary residue defects.

\subsection*{D-branes}
Chan--Paton multiplicity forces ordered multiplication and therefore an $E_1$
rather than an $E_\infty$ structure.  The nonabelianity of the first planted-forest
higher bracket is the local residue shadow of that ordered brane geometry.

\subsection*{Holography}
The nested commutator
\[
[\beta,[\beta,\chi(\alpha)]]
\]
is the first local Yangian-like pattern in the story: a boundary mismatch acted
on twice by an overlap parameter.  It is the smallest local seed of the larger
transport from modular boundary data to line-operator algebra.

\section{What this changes}

The logic of the subject becomes cleaner.

\begin{enumerate}[label=\textup{(\arabic*)}]
\item The local renormalization problem is solved once the positive normal-weight
  ideal is identified.  The local singular sector is contractible.
\item The surviving local deformation theory is residue-theoretic.
\item The first nontrivial higher operation occurs on $\FM_3(X)$ and is the
  nested-collision bracket $l_3$.
\item The first genuinely higher local obstruction is the cubic planted-forest
  obstruction of Theorem~\ref{thm:fm3-higher-obstruction-synthesis}.
\item The next frontier is no longer local contraction.  It is global
  higher-depth gluing on the full planted-forest stratification of $\FM_n(X)$
  and its higher-genus analogues.
\end{enumerate}

In that precise sense, the $\FM_3$ model is the platonic first atom of modular
homotopy theory beyond strict dg Lie geometry: it is the smallest space on which
local residues, higher operations, and the geometry of nested collisions become
inseparable.

